\documentclass{ximera}

\title{Ladingen}
\author{Daan Lipkens}

\begin{document}
\begin{abstract}
    Ladingshoeveelheid en en soorten ladingen.
\end{abstract}
\maketitle

Er zijn simpele experimentjes om aan te tonen dat er elektrische krachten zijn. Zo gaan je haren rechtstaan als je met een ballon over je haar wrijft. Of het aantrekken van papieren snippers met een wollen trui waarover je hebt gewreven.
%--------------------------
%yt video naar demonstratie met papier snippers --> link blijft error geven met \youtube
    %\youtube{https://youtu.be/Glifj-apltU?si=gRhOZ2T3wFJypikX}
	\mylink{https://www.youtube.com/watch?v=Glifj-apltU}{Link naar video met het aantrekken van papiersnippers.}
%--------------------------


% -------------------------
%	eigenschappen van ladingen en soorten ladingen
% -------------------------------

In een reeks van simpele experimenten, zoals te zien op figuur \ref{fig:buisjes}, kan er aangetoond worden dat er maar twee soorten ladingen bestaan en is als gevolg hiervan het deeltjes model uitgebreid:
%--------------------------
%afbeelding buisjes --> verschijnen niet in html
\begin{figure}[h]
\captionsetup{width=0.85\textwidth}
\begin{center}
\includegraphics[width=0.8\textwidth]{/pic5/elektriciteit/coulomb/buisjes.jpeg}
\caption{(a) Een negatief geladen rubberen staaf die aan een draad hangt, wordt aangetrokken door een positief geladen glazen staaf. (b) Een negatief geladen rubberen staaf wordt afgestoten door een andere negatief geladen rubberen staaf.}
\label{fig:buisjes}
\end{center}
\end{figure}
%--------------------------
\begin{itemize}
 	\item er bestaan twee soorten ladingen: \textcolor{red}{\textbf{positieve}} ladingen (+) en \textcolor{blue}{\textbf{negatieve}} ladingen (-). \footnote{Conventie voor de kleuren: \textcolor{red}{rood} voor \textcolor{red}{positieve} lading en \textcolor{blue}{blauw}/zwart voor \textcolor{blue}{negatieve} ladingen}
	\item de kracht dat ladingen op elkaar uitoefenen is de elektrische kracht. \textbf{Gelijksoortige ladingen stoten elkaar af} en \textbf{verschillende soorten ladingen trekken elkaar aan}.
	\item \textbf{lading blijft behouden}. Lading kan niet vernietigd of gemaakt worden, enkel uitgewisseld.
	\end{itemize}
	\begin{proposition}
	\textbf{Behoud van lading}\\ De totale elektrische lading in een gesloten systeem blijft altijd gelijk. Met andere woorden: lading kan niet uit het niets ontstaan of zomaar verdwijnen. Ze kan enkel worden verplaatst van het ene voorwerp naar het andere. Wanneer een voorwerp elektronen verliest, neemt een ander voorwerp ze op, zodat de totale hoeveelheid lading onveranderd blijft.
	\end{proposition}
	\begin{itemize}
	\item lading is \textbf{gekwantiseerd}\footnote{Gekwantiseerd betekent dat een bepaalde grootheid niet willekeurig elke waarde kan aannemen, maar enkel specifieke, discrete waarden. Deze waarden wordt \textit{e} genoemd. $\mathit{e}=1,6\cdot 10^{-19}\unit{\coulomb}$. De huidige theorieën laten wel breuken van deze lading toe voor quarks, al zijn \textit{vrije} quarks nog nooit gevonden. Meer hier over in \ref{hoofdstuk: Deeltjes Fysica}}%hoofdstuk~\ref{hoofdstuk: Deeltjes Fysica} 
\end{itemize}
\begin{definition}
De \textbf{ladingshoeveelheid}, of ook wel elektrische lading, stellen we voor door \textbf{Q}. De eenheid van ladingshoeveelheid is de coulomb (C).
\end{definition}


Lange tijd dacht men dat lading en materie twee gescheiden domeinen van de wetenschap waren. Onderzoek heeft echter aangetoond dat lading voorkomt in de deeltjes waaruit ook materie is opgebouwd. Materie bestaat uit atomen, en atomen bestaan uit protonen, neutronen en elektronen. Meer informatie over de bouwstenen van alle deeltjes volgt in een later hoofdstuk: hoofdstuk~\ref{hoofdstuk: De Atoomkern}.

% TABEL MET GEGEVENS NUCLEONEN
\begin{table}[h!]
\caption{Lading en Massa van de Proton, Neutron en het Elektron}
\centering
\begin{tabular}{@{}lcr@{}}
	\toprule
	deeltje & lading (C) & massa (kg)\\
	\midrule
	\textcolor{red}{proton} & $+1,60\cdot 10^{-19}$ & $1,673\cdot 10^{-27}$ \\
	neutron & 0 & $1,675\cdot 10^{-27}$ \\
	\textcolor{blue}{elektron} & $-1,60\cdot 10^{-19}$ & $9,109\cdot 10^{-31}$ \\
	\bottomrule
\end{tabular}
\vspace{2mm}
	\begin{remark}
		De ladingshoeveelheid is vanbuiten te kennen, de massa van de deeltjes niet. Coulomb is een grote eenheid, dus vaak zal er gebruik worden gemaakt van voorvoegsels zoals mili, micro en nano. Zie \ref{hoofdstuk: Standaardvorm}.%appendix~\ref{hoofdstuk: Standaardvorm}
	\end{remark}
\end{table}


\end{document}
